\documentclass{article}
\title{Practice Problem 2.39}
\author{CSAPP}
\usepackage{booktabs, parskip}

\begin{document}
    \maketitle
    \begin{flushleft}
       For a run of ones starting at bit position $n$ down to bit position
       $m(n\geq m)$, we saw that we can generate two forms of code, $A$ and $B$.
       How should the compiler decide which form to use?

       \textbf{\small Solution}\\
        For a run of $k = n - m + 1$ ones, Form A requires $k$ shifts and
    $k-1$ additions, so its cost grows linearly with $k$ (roughly $2k-1$
    instructions). Form B requires at most $2$ shifts and $1$ subtraction,
    a constant cost regardless of $k$. Comparing the two:
    \begin{itemize}
        \item $k = 1$: Form A costs $1$ instruction; Form B costs $2$--$3$.
            Form A wins.
        \item $k = 2$: Form A costs $3$ instructions; Form B costs $2$--$3$.
            Roughly tied.
        \item $k \geq 3$: Form A's cost ($2k-1$) exceeds Form B's constant
            cost. Form B wins.
    \end{itemize}
    The compiler should therefore use Form A for short runs ($k \leq 2$)
    and switch to Form B once the run length reaches $k = 3$ or more,
    since at that point the savings from collapsing the run outweigh the
    fixed overhead of the extra subtraction.
       
    \end{flushleft}
\end{document}